
\chapter{Introducción}

\section{Contexto}

El desarrollo de vehículos de competición modernos requiere
la integración de múltiples disciplinas de ingeniería,
matemáticas aplicadas y ciencias computacionales.

El diseño de un vehículo de alto rendimiento involucra
problemas relacionados con la aerodinámica, la dinámica
vehicular, la selección de materiales, la optimización
geométrica y el análisis de sistemas físicos complejos.

La resolución de estos problemas puede requerir una cantidad
considerable de recursos computacionales, especialmente
cuando se utilizan métodos numéricos, simulaciones y procesos
iterativos de optimización.

En este contexto, las herramientas de computación científica
y las técnicas de inteligencia artificial ofrecen
oportunidades para desarrollar sistemas que faciliten
el análisis y la exploración de soluciones de ingeniería.

\section{Motivación}

APEX surge como una iniciativa de desarrollo e investigación
orientada a explorar la aplicación de las ciencias
computacionales a problemas de ingeniería de vehículos
de competición.

El proyecto busca integrar infraestructura computacional,
modelos matemáticos, herramientas de simulación y técnicas
de optimización dentro de una plataforma modular.

La motivación principal es investigar cómo estas tecnologías
pueden combinarse para construir herramientas capaces de
asistir en procesos de diseño, análisis y evaluación
de sistemas de ingeniería.

\section{Objetivo general}

Desarrollar progresivamente una plataforma de ingeniería
computacional que integre herramientas de procesamiento,
modelado, simulación y optimización para el análisis
y diseño de vehículos de alto rendimiento.

\section{Objetivos específicos}

\begin{itemize}
    \item Diseñar una arquitectura de software modular
    y extensible.

    \item Implementar mecanismos de detección y gestión
    de recursos computacionales.

    \item Desarrollar infraestructura para la ejecución
    paralela de tareas.

    \item Establecer mecanismos de registro y trazabilidad
    de las operaciones del sistema.

    \item Investigar e integrar modelos matemáticos
    y físicos aplicados a la ingeniería vehicular.

    \item Evaluar experimentalmente los componentes
    desarrollados mediante pruebas reproducibles.
\end{itemize}

\section{Alcance del documento}

Este documento registra el desarrollo técnico y científico
de APEX. Describe sus fundamentos, decisiones de arquitectura,
métodos de implementación y procedimientos de evaluación.

Se distinguirán explícitamente los componentes implementados,
los métodos en investigación y las capacidades pendientes
de desarrollo y validación.

En su estado actual, el proyecto se encuentra en una etapa
inicial centrada en la infraestructura computacional.
Por tanto, las capacidades de simulación aerodinámica,
optimización física y diseño autónomo se consideran
objetivos de investigación futuros y no funcionalidades
actualmente demostradas.
